\section{A transfer in the maximal Roe algebra}\label{sec:transfer}

The projection $f:Y\to B$ forgets a finite Cayley graph over each base
vertex.  We will lift operators on $B$ to operators on $Y$ by transporting
their matrix entries between these fibers and then projecting onto the
constant vectors.  The resulting map supplies the analytic class that we
will move from one end of $X$ to the other.

\begin{proposition}\label{ana:transfer}
For the projection $f:Y\to B$ there is a $*$-homomorphism
\[
 \Phi:A_B\longrightarrow A_Y
\]
such that $(f_A\Phi)_*=\id$ on $K_*(A_B)$.  In particular,
\[
 (f_A)_*:K_1(A_Y)\longrightarrow K_1(A_B)
\]
is surjective.
\end{proposition}

Throughout this section, $u,v,w$ denote vertices of $B$.  Write
$F_u=f^{-1}(u)$ and let
\[
 \xi_u=|F_u|^{-1/2}\sum_{y\in F_u}\delta_y
 \quad\text{in }\ell^2(F_u).
\]
These vectors define an isometry
\begin{equation}\label{eq:constant-isometry}
 V:H_B\longrightarrow H_Y,
 \qquad V(\delta_u\otimes h)=\xi_u\otimes h.
\end{equation}
The diameters of the fibers are not uniformly bounded, so $V$ itself need
not be a covering isometry for a coarse map.  We will obtain an element of
the maximal Roe algebra from this averaging operation by using
property~\textup{(T)}.

\subsection{The projection onto fiber constants}

For $\gamma\in\Gamma$, let $U_\gamma$ act by left multiplication by
$\gamma$ on every fiber, with the identity on $H_0$.  This is a unitary
operator of propagation at most the word length of $\gamma$.
Lemma~\ref{ana:multipliers} therefore gives a unitary representation
\[
 \Gamma\longrightarrow M(A_Y),\qquad\gamma\longmapsto U_\gamma,
\]
and hence a unital $*$-homomorphism from $C^*_{\max}(\Gamma)$ into
$M(A_Y)$.

We recall explicitly the consequence of property~\textup{(T)} that will
be used; see \cite{BHV08}.  In $C^*_{\max}(\Gamma)$ put
\[
 \Delta=1-\frac1{|S|}\sum_{s\in S}s.
\]
For every unitary representation $\pi$ and every vector $\eta$,
\[
 \langle\pi(\Delta)\eta,\eta\rangle
   =\frac1{2|S|}\sum_{s\in S}
      \|\pi(s)\eta-\eta\|^2.
\]
In particular, $0\le\Delta\le2$, and the kernel of $\pi(\Delta)$ is
the invariant subspace.  Property~\textup{(T)}, applied to the fixed finite
generating set $S$, gives $\kappa>0$ such that every unit vector
orthogonal to that subspace is moved by at least $\kappa$ by some
$s\in S$.  Consequently, with $\varepsilon=\kappa^2/(2|S|)$,
\[
 \sigma_{C^*_{\max}(\Gamma)}(\Delta)
     \subseteq\{0\}\cup[\varepsilon,2].
\]
The continuous functional calculus thus gives a projection onto the
invariant vectors in every representation, and
\[
 \left(1-\frac{\Delta}{2}\right)^n
\]
converges to this projection in the maximal group norm.  Denote its image
in $M(A_Y)$ by $P$, and denote the images of the displayed approximants
by $P_n$.  We have
\begin{equation}\label{eq:kazhdan-convergence}
 \|P_n-P\|_{M(A_Y)}\longrightarrow0.
\end{equation}
Each $P_n$ is a finite-propagation multiplier preserving every fiber.

In the concrete representation on $H_Y$, the invariant vectors in
$\ell^2(F_u)$ are precisely the constant vectors: the map
$\Gamma\to Q_{2j}$ is onto and its left action on $F_u$ is transitive.
Thus the concrete image of $P$ is $VV^*$. The transfer itself will use
$P$ in the maximal multiplier algebra $M(A_Y)$.

\subsection{Transporting matrix entries}

For each pair $u,v$ in the same component $B_j$, choose a shortest edge
path from $v$ to $u$, choosing the reverse path for the reversed pair.
Horizontal transport along this path is a permutation
\[
 O_{uv}:\ell^2(F_v)\longrightarrow\ell^2(F_u).
\]
Take $O_{uu}=1$.  When $u,v$ lie in different components, instead put
\[
 O_{uv}=|\xi_u\rangle\langle\xi_v|.
\]
In both cases,
\begin{equation}\label{eq:transport-constants}
 O_{vu}=O_{uv}^*,\qquad
 O_{uv}\xi_v=\xi_u,\qquad O_{uv}^*\xi_u=\xi_v.
\end{equation}
Moreover, $O_{uv}$ intertwines the left $\Gamma$-actions.  For
horizontal transport this follows because the matchings use right
multiplication; for the rank-one operators it follows because both
vectors are invariant.

Given $T\in\C[B]$, use the decomposition
$H_Y=\bigoplus_{u\in B}\ell^2(F_u)\otimes H_0$ to define
\begin{equation}\label{eq:transport-lift}
 L(T)_{uv}=O_{uv}\otimes T_{uv}.
\end{equation}
This is a linear, adjoint-preserving lift, but it need not be
multiplicative.  The next lemma isolates both its control and the precise
failure of multiplicativity.

\begin{lemma}\label{ana:transport-defect}
For $T\in\C[B]$, the operator $L(T)$ belongs to $\C[Y]$ and commutes
with $P$.  For $T,T'\in\C[B]$, the defect
\begin{equation}\label{eq:transport-defect}
 D=L(T)L(T')-L(TT')
\end{equation}
has finite vertex support and satisfies $PD=0$ in $A_Y$.
\end{lemma}

\begin{proof}
Let $R$ be a propagation bound for $T$.  The matrix support of $T$ can
be partitioned into finitely many matchings of the base vertices, with a
number depending only on $R$ and the bounded geometry of $B$.  On one
such matching, the blocks in \eqref{eq:transport-lift} have disjoint
domain and range blocks, and their norms are at most $\|T\|_{\mathrm{op}}$.
Summing over the matchings proves that $L(T)$ is bounded.

Within a component, any nonzero vertex entry of $O_{uv}$ joins vertices
connected by the chosen horizontal path.  Its length is
$d_B(u,v)\le R$.  Among different components, only finitely many vertex
pairs $(u,v)$ can occur at this propagation bound; their fibers are
finite.  The remaining entries therefore also have a finite propagation
bound.  Each vertex entry of $L(T)$ is a scalar multiple of a compact
operator $T_{uv}$.  It follows that $L(T)\in\C[Y]$.
The intertwining property of $O_{uv}$ shows that $L(T)$ commutes with
every $U_\gamma$, hence with every $P_n$ and with $P$.

We next locate the defect.  Let $R'$ be a propagation bound for $T'$.
Outside finitely many components, all entries involved in
\eqref{eq:transport-defect} stay in a single component $B_j$.
For a summand with intermediate vertex $v$, the concatenated path from
$w$ to $v$ to $u$ has length at most $R+R'$, as does the chosen path
from $w$ to $u$.  If $l_j>2(R+R')$, these two paths have the same total
seam crossing numbers.  Indeed, their lifted coordinate displacements
differ by a multiple of $l_j$ and have difference of absolute value at
most $2(R+R')$ in each coordinate.  That multiple must be zero.
Since $U_j$ and $V_j$ commute in $Q_{2j}$, horizontal transport depends
only on these seam totals.  It follows that
\[
 O_{uv}O_{vw}=O_{uw}
\]
for every summand under consideration.  The defect therefore vanishes
on all sufficiently large components.  The omitted components and any
cross-component entries lie in a finite union of finite graphs, so $D$
has finite vertex support.

Equations~\eqref{eq:transport-constants} give, on the concrete Hilbert
spaces,
\[
 L(T)V=VT,\qquad V^*L(T)=TV^*.
\]
The concrete image of $PD$ is consequently zero.  To obtain the equality
in the maximal completion, choose a finite union $E$ of components of
$Y$ containing the vertex support of $D$.  Every $P_n$ preserves each
fiber, so $P_nD$ has vertex support in the same finite set $E$.
Lemma~\ref{ana:compacts} identifies its maximal norm with its concrete
operator norm.  On $H_Y$, we have $P_nD\to0$ in operator norm; hence
$P_nD\to0$ in maximal norm.  On the other hand,
\eqref{eq:kazhdan-convergence} gives $P_nD\to PD$ in $A_Y$.
Thus $PD=0$ in $A_Y$.
\end{proof}

\subsection{The splitting}

We can now define the transfer and compute its composition with $f_A$.
The first task used finite vertex support.  The second will use a
different control: averaging can increase propagation inside a fiber,
but it does not increase propagation after the fiber is projected to its
base vertex.

\begin{proof}[Proof of Proposition~\ref{ana:transfer}]
For $T\in\C[B]$, put
\[
 \Phi(T)=PL(T)\in A_Y.
\]
Since $P$ is a projection commuting with the lifts,
Lemma~\ref{ana:transport-defect} gives
\[
 \Phi(T)\Phi(T')=P L(T)L(T')=P L(TT')=\Phi(TT').
\]
The lift preserves adjoints, so $\Phi$ is a $*$-homomorphism on
$\C[B]$.  The universal property of the maximal norm extends it to a
$*$-homomorphism $A_B\to A_Y$.

Let $J:H_Y\to H_B$ be the covering isometry defining $f_A$; it sends
the summand at $y\in F_u$ isometrically into the summand at $u$, with
orthogonal ranges for distinct $y$ in the same fiber.  For
$T\in\C[B]$, continuity gives
\begin{equation}\label{eq:split-approximants}
 f_A\Phi(T)=\lim_{n\to\infty} J P_nL(T)J^*
 \qquad\text{in }A_B.
\end{equation}
All operators in this limit have propagation at most $\prop(T)$ on
$B$, because $P_n$ preserves the fibers.  In concrete operator norm
they converge to
\[
 JVV^*L(T)J^*=JVT V^*J^*.
\]
The isometry $JV:H_B\to H_B$ lies over the identity of $B$, so
$JVT V^*J^*$ belongs to $\C[B]$ and has the same propagation bound.
The difference from each approximant in
\eqref{eq:split-approximants} has this fixed propagation bound as well.
Lemma~\ref{ana:fixed-propagation} promotes the concrete convergence to
maximal-norm convergence.  Hence
\begin{equation}\label{eq:split-identity}
 f_A\Phi(T)=\Ad(JV)(T).
\end{equation}
By density this identity holds on $A_B$.
Finally, Proposition~\ref{ana:functoriality} says that a covering isometry
over the identity induces the identity on $K$-theory.  Applying it to
$JV$ proves $(f_A\Phi)_*=\id$.
\end{proof}
